% concatenated from besz10v2.tex
\documentclass [12pt]{article}
\usepackage{amssymb,amsmath,amsthm}
\usepackage{color}
\def\bit{\bfseries\textit}
\def\bif{\textbf}
\def\Im{{\rm Im}\,}
\def\tg{{\,\rm tg}\,}
\DeclareMathOperator{\spn}{span}
\DeclareMathOperator{\ind}{ind}
\DeclareMathOperator{\tr}{tr}
\DeclareMathOperator{\dom}{dom}
\def\arctg{{\rm arctg\,}}
\usepackage[cp1250]{inputenc}
\usepackage[active]{srcltx}
\def\d{\displaystyle}
\def\t#1{$^{\bf *}$\bf #1{.}}
\newtheorem{thm}{Theorem}[section]
\newtheorem{prop}[thm]{Proposition}
\newtheorem{lem}[thm]{Lemma}
 \newtheorem{rem}[thm]{Remark}
 \newtheorem{cor}[thm]{Corollary}
\newtheorem{defn}[thm]{Definition}
\newtheorem{ex}[thm]{Example}
\DeclareMathOperator{\supp}{supp}
\DeclareMathOperator{\ext}{ext}
\newcommand{\C}{\mathbb{C}}
\newcommand{\R}{\mathbb{R}}
\renewcommand{\S}{\mathbb{S}}
\newcommand{\N}{\mathbb{N}}
\newcommand{\T}{\mathbb{T}}
\newcommand{\M}{\mathbb{M}}
\newcommand{\D}{\mathbb{D}}
\newcommand{\Z}{\mathbb{Z}}
\newcommand{\la}{\lambda}
\def\labelenumi{\bf \arabic{enumi}.}
\def\tlab{$\bf *$\labelenumi}
\long\def\comment#1{{}}
\def\ph{\phantom{-}}
\def\p{\ph}
 \title{Indeterminate Jacobi operators II}
 \author{Christian Berg and Ryszard Szwarc}
 

\begin{document}

 \maketitle
file: besz10v2.tex

\begin{abstract}  We consider the Jacobi operator $(T,D(T))$ associated with an indeterminate Hamburger moment problem, and present countable subsets $S$ of the domain $D(T)$  such that $\spn(S)$ is dense in $\ell^2$. As an example  we have $S=\{(p_n(u))+B(u)(p_n(0)) \mid D(u)=0, u\neq 0\}$, where $(p_n)$ denotes the orthonormal polynomials of the moment problem and $B, D$ are two of the Nevanlinna functions. 
It is also proved that sets like $S$ are optimal in the sense that if one vector is removed, then the span is no longer dense.
\end{abstract}

{\bf Mathematics Subject Classification}: Primary 47B25; Secondary 47B36, 44A60

{\bf Keywords}. Jacobi matrices, domains of Jacobi operators, indeterminate moment problems.


\section{Introduction and main results} This paper is a continuation of \cite{B:S5}  by improving certain results about the domain of the indeterminate Jacobi operator. We  consider the Jacobi matrix $J$ associated with a moment sequence $s=(s_n)_{n\ge 0}$ of the form
\begin{equation}\label{eq:mom}
s_n=\int x^n\,d\mu(x),\quad n=0,1,\ldots,
\end{equation}
where $\mu$ is a positive measure on $\R$  with infinite support and moments of every order. It is a tridiagonal matrix of the  form
\begin{equation}\label{eq:Jac}
J=\begin{pmatrix}
b_0 & a_0 & 0 & \hdots\\
a_0 & b_1 & a_1 & \hdots\\
0 & a_1 & b_2 & \hdots\\
\vdots &\vdots & \vdots & \ddots
\end{pmatrix},
\end{equation}
where $a_n>0, b_n\in\R, n\ge 0$  are given by the three term recurrence relation
\begin{equation*}\label{eq:3term}
xp_n(x)=a_np_{n+1}(x)+b_np_n(x)+a_{n-1}p_{n-1}(x),\quad n\ge 0, \quad a_{-1}:=0.
\end{equation*}
Here $(p_n)_{n\ge 0}$ is the sequence of orthonormal polynomials associated with $\mu$, hence satisfying
\begin{equation*}\label{eq:op}
\int p_n(x)p_m(x)\,d\mu(x)=\delta_{n,m}.
\end{equation*}
As usual $p_n$ is a real polynomial of degree $n$ with positive leading coefficient. As in \cite{B:S5} we follow the terminology of \cite{Sch}. Basic results about the classical moment problem can also be found in \cite{Ak} and \cite{S}. Recent  results about indeterminate moment problems can be found in \cite{B:S1}, \cite{B:S2}, \cite{B:S3}, \cite{B:S4}, \cite{B:S5}.

 It is easy to see that the proportional measures $\la\mu, \la>0$ lead to the same Jacobi matrix $J$, and the well-known Theorem of Favard (see \cite[Theorem 5.14]{Sch}) states that any matrix of the form \eqref{eq:Jac} with $a_n>0, b_n\in\R$ comes from a unique  moment sequence  $(s_n)$ as above, normalized such that $s_0=1$. In the following we shall always assume that this normalization holds, and consequently the solutions $\mu$ of \eqref{eq:mom} are probability measures  and $p_0=1$.

The Jacobi matrix acts as a symmetric operator  in the Hilbert space $\ell^2$ of square summable complex sequences. Its domain $\mathcal F$ consists of the complex sequences $(c_n)_{n\ge 0}$ with only finitely many non-zero terms, and the action is multiplication of the matrix $J$ by $c\in\mathcal F$ considered as a column, i.e.,
\begin{equation}\label{eq:J} 
(Jc)_n:=a_{n-1}c_{n-1}+b_n c_n+a_n c_{n+1},\quad n\ge 0.
\end{equation}
Denoting $(e_n)_{n\ge 0}$ the standard orthonormal basis of $\ell^2$, we have
\begin{equation*}\label{eq:F}
\mathcal F=\spn \{e_n|n\ge 0\}.
\end{equation*}

\begin{defn} The Jacobi operator associated with $J$ is by definition the closure $(T,D(T))$ of the symmetric operator $(J,\mathcal F)$.
\end{defn}

It is a classical fact that the closed symmetric operator $(T,D(T))$ has deficiency indices  either $(0,0)$  or $(1,1)$. These cases occur precisely if the  moment sequence \eqref{eq:mom} is {\it determinate} or {\it indeterminate}, i.e., there is exactly one or several solutions $\mu$ satisfying \eqref{eq:mom}.   

 By definition $D(T)$ consists of those $c\in\ell^2$ for which there exists a sequence
 $(c^{(k)})\in\mathcal F$ such that $\lim_{k\to\infty}c^{(k)}=c$ and $(Jc^{(k)})$ is a convergent sequence in $\ell^2$. For  such $c$ we have $Tc=\lim_{k\to\infty} Jc^{(k)}$, and this limit is independent of the choice of approximating sequence $(c^{(k)})$. 
 
 Clearly, $D(T)$ is closed under complex conjugation and 
 $$
 T\overline{c}=\overline{Tc},\quad c\in D(T).
 $$
 
 We recall that the adjoint operator $(T^*,D(T^*))$ is the maximal operator associated with $J$, cf. \cite[Proposition 6.5]{Sch}. In fact, the matrix product of $J$  and any column vector $c$ makes sense, cf. \eqref{eq:J}, and $D(T^*)$ consists of those $c\in\ell^2$ for which the product $Jc$ belongs to $\ell^2$. For $c\in D(T^*)$ we have $T^*c=Jc$.
 
In this paper we will only consider the  indeterminate case of the Jacobi operator $(T,D(T))$, where it is known that  the set of solutions $\mu$ to \eqref{eq:mom} is an infinite convex set $V$. The polynomials of the second kind $(q_n)$ are given as
\begin{equation*}
q_n(z)=\int \frac{p_n(z)-p_n(x)}{z-x}\,d\mu(x), \quad z\in \C,
\end{equation*}
where $\mu\in V$ is arbitrary. Note that $q_n$ as well as $p_n$ are independent of $\mu\in V$ in the indeterminate case.

We define  the sequences
\begin{equation}\label{eq:frak}
 \mathfrak{p}_z:=(p_n(z)), \mathfrak{q}_z:=(q_n(z)),\quad z\in\C,
 \end{equation}
 where we have followed the terminology of \cite{Sch}. It is known that they belong to
$\ell^2$ because of indeterminacy, and
   $||\mathfrak{p}_z||$ and $||\mathfrak{q}_z||$ are positive continuous functions for $z\in\C$.  It is therefore possible 
for $c\in\ell^2$ to define entire functions $F_c, G_c$ as
\begin{equation}\label{eq:FGc}
F_c(z)=\sum_{n=0}^\infty c_np_n(z),\quad G_c(z)=\sum_{n=0}^\infty c_nq_n(z),\quad z\in\C.
\end{equation}

It is well known that $F_c\in L^2(\mu)$ for any solution $\mu\in V$ and that
$$
\lim_{n\to\infty} \sum_{k=0}^n c_kp_k(z)= F_c(z)
$$
locally uniformly in $z\in\C$ and in $L^2(\mu)$ for any $\mu\in V$. Furthermore, Parseval's equation holds
\begin{equation}\label{eq:Pars}
\int |F_c(x)|^2\,d\mu(x)=||c||^2,\quad c\in\ell^2, \mu\in V.
\end{equation}

We recall the following four entire functions of two complex variables, called the {\it Nevanlinna functions} of the indeterminate moment problem:
\begin{eqnarray}
A(u,v)&=&(u-v)\sum_{k=0}^\infty q_k(u)q_k(v)\label{eq:A},\\
B(u,v)&=&-1+(u-v)\sum_{k=0}^\infty p_k(u)q_k(v) \label{eq:B},\\
C(u,v)&=&1+(u-v)\sum_{k=0}^\infty q_k(u)p_k(v)\label{eq:C},\\
D(u,v)&=&(u-v)\sum_{k=0}^\infty p_k(u)p_k(v)\label{eq:D},
\end{eqnarray}
see Section 7.1 in \cite{Sch}. They satisfy the  fundamental determinant equation
\begin{equation}\label{eq:2vardet}
A(u,v)D(u,v)-B(u,v)C(u,v)=1, \quad u,v\in\C.
\end{equation}
 
We define entire functions of one variable by setting the second variable to 0, i.e., 
\begin{equation}\label{eq:A-D}
A(u)=A(u,0),\;B(u)=B(u,0),\;C(u)=C(u,0),\;D(u)=D(u,0),
\end{equation}
and \eqref{eq:2vardet} becomes

\begin{equation}\label{eq:det}
A(u)D(u)-B(u)C(u)=1,\quad u\in\C.
\end{equation}

By Section 6.5 in \cite{Sch} we have
\begin{equation}\label{eq:p,q}
\mathfrak{p}_z, \mathfrak{q}_z \in D(T^*),\quad T^*\mathfrak{p}_z=z\mathfrak{p}_z, 
\quad T^*\mathfrak{q}_z=e_0+z\mathfrak{q}_z,\quad z\in\C.
\end{equation}

A main result of \cite{B:S5} states the following:

\begin{thm}\label{thm:pqmain} For all $z\in\C$ we have $\mathfrak{p}_z, \mathfrak{q}_z\notin D(T)$.

Let $u,v\in\C$ be given.
\begin{enumerate}
\item[(i)] There exists $\alpha\in\C$ such that $\mathfrak{p}_u+\alpha \mathfrak{p}_v\in D(T)$ if and only if $D(u,v)=0$. In the affirmative case $\alpha$  is uniquely determined as
$\alpha=B(u,v)$. 

\item[(ii)] There exists $\beta\in\C$ such that $\mathfrak{q}_u+\beta \mathfrak{q}_v\in D(T)$ if and only if $A(u,v)=0$. In the affirmative case $\beta$  is uniquely determined as $\beta=-C(u,v)$.

\item[(iii)] There exists $\gamma\in\C$ such that $\mathfrak{p}_u+\gamma \mathfrak{q}_v\in D(T)$ if and only if $B(u,v)=0$. In the affirmative case $\gamma$  is uniquely determined as $\gamma=-D(u,v)$. In particular $\mathfrak{p}_u+\gamma\mathfrak{q}_u\notin D(T)$ for all $u, \gamma\in\C$.
\end{enumerate}
\end{thm}

From  this theorem we have the following concrete subspaces of $D(T)$:
\begin{eqnarray*}
P&=&\spn\{\mathfrak{p}_u+B(u,v){\mathfrak{p}}_v \mid u,v\in\C, D(u,v)=0, u\neq v\},\\
Q&=&\spn\{\mathfrak{q}_u-C(u,v){\mathfrak{q}}_v \mid u,v\in\C, A(u,v)=0, u\neq v\},\\
M&=&\spn\{\mathfrak{p}_u-D(u,v){\mathfrak{q}}_v \mid u,v\in\C, B(u,v)=0\}.
\end{eqnarray*}
Note that $\mathfrak{p}_u+B(u,v)\mathfrak{p}_{v} =\mathfrak{q}_u-C(u,v)\mathfrak{q}_v=0$  for $u=v$ and that $Q\subseteq e_0^\perp=\{(c_n)_{n\ge 0}\in\ell^2 \mid c_0=0\}$.

For a fixed number $v_0\in\R$ we define the following subspaces of $P,Q,M$
respectively
\begin{eqnarray}
P(v_0)&=&\spn\{\mathfrak{p}_u+B(u,v_0)\mathfrak{p}_{v_0} \mid  D(u,v_0)=0, u\neq v_0\},\label{eq:P}\\
Q(v_0)&=&\spn\{\mathfrak{q}_u-C(u,v_0){\mathfrak{q}}_{v_0} \mid  A(u,v_0)=0, u\neq v_0\},\label{eq:Q}\\
M(v_0)&=&\spn\{\mathfrak{p}_u-D(u,v_0){\mathfrak{q}}_{v_0} \mid  B(u,v_0)=0\}\label{eq:M}.
\end{eqnarray}
In these definitions it is important to remember that if $F$ is any of the functions $A,B,C,D$ of two variables, then $\{u\in\C \mid F(u,v_0)=0\}$
is a countably infinite set of real numbers, cf. Theorem 1.3 in \cite{B:S5}.

Letting $\mu[v_0]$ denote the unique N-extremal measure in $V$ with $v_0\in\supp(\mu[v_0])$,
cf. Proposition~\ref{thm:supportNext}, we have by Theorem 3 in \cite{B:C}
\begin{equation}\label{eq:supp}
\supp(\mu[v_0])=\{u\in\R \mid D(u,v_0)=0\},
\end{equation}
and hence
\begin{equation} 
P(v_0)=\spn\{\mathfrak{p}_u+B(u,v_0)\mathfrak{p}_{v_0} \mid u\in\supp(\mu[v_0])
\setminus\{v_0\}\}.
\end{equation}

In the next section we recall the parametrization $\mu_t, t\in\R^*=\R\cup\{\infty\}$ of the N-extremal measures in $V$, and by Proposition~\ref{thm:supportNext} we have
\begin{equation}\label{eq:N-ext}
\mu[v_0]=\mu_t,\quad t=-B(v_0)/D(v_0)
\end{equation}
with the convention that $t=\infty$ if $D(v_0)=0$.

Also the set $\{u\in\R \mid B(u,v_0)=0\}$ is the support of an N-extremal measure. By formula (5.6) in \cite{B:S5} we have $B(u,v_0)=B(u)C(v_0)-D(u)A(v_0)$
so
\begin{equation}\label{eq:N-ext2}
\{u\in\R \mid B(u,v_0)=0\}=\supp(\mu_s),
\end{equation}
where $s=-A(v_0)/C(v_0)$ interpreted as $s=\infty$ if $C(v_0)=0$.

 The following result sharpens that $D(T)$ is dense in $\ell^2$ by giving concrete dense subspaces of $D(T)$, which are all spanned by countably many vectors. It is the main result of this paper, but we believe that the results of Proposition~\ref{thm:1-2} and Corollary~\ref{thm:limit} are new.

  \begin{thm}\label{thm:concrete} For each $v_0\in\R$ the subspaces $P(v_0)$ and $M(v_0)$ are dense in $\ell^2$, while $Q(v_0)$ is dense in $e_0^\perp$.
\end{thm}

\begin{rem}\label{thm:optimal1}
{\rm We claim that the subspaces $P(v_0), v_0\in \R$ are optimal as dense subspaces in $\ell^2$ in the following sense. Let $D(u_2,v_0)=0$ for $u_2\neq v_0$ and define
\begin{equation}\label{eq:P'} 
P'(v_0)=\spn\{\mathfrak{p}_u+B(u,v_0)\mathfrak{p}_{v_0} \mid D(u,v_0)=0, u\neq v_0,u_2\}.
\end{equation}
Then $P'(v_0)$ is not dense in $\ell^2$.

Similarly $Q(v_0), v_0\in\R$ are optimal as  dense subspaces of $e_0^\perp$ in the sense that if $A(u_2,v_0)=0, u_2\neq v_0$ then
\begin{equation}\label{eq:Q'} 
Q'(v_0)=\spn\{\mathfrak{q}_u-C(u,v_0)\mathfrak{q}_{v_0} \mid A(u,v_0)=0, u\neq v_0,u_2\}
\end{equation}
is not dense in $e_0^\perp$.
}
\end{rem}

\begin{rem}\label{thm:optimal2}
{\rm Let $u_1,u_2$ be different real numbers satisfying $B(u_1,v_0)=B(u_2,v_0)=0$ and define
\begin{eqnarray} 
M'(v_0)&=&\spn\{\mathfrak{p}_u-D(u,v_0)\mathfrak{q}_{v_0} \mid B(u,v_0)=0, u\neq u_1\},\label{eq:M'}\\
M''(v_0)&=&\spn\{\mathfrak{p}_u-D(u,v_0)\mathfrak{q}_{v_0} \mid B(u,v_0)=0, u\neq u_1, u_2\},\label{eq:M''}
\end{eqnarray}
then $M'(v_0)$ is dense in $\ell^2$ while $M''(v_0)$ is not dense in $\ell^2$.
}
\end{rem}


The proof of Theorem~\ref{thm:concrete} and  Remarks~\ref{thm:optimal1} and \ref{thm:optimal2}  will be given in Section 3.





\section{Preliminaries about indeterminate moment problems}

For the proof of Theorem~\ref{thm:concrete} we need the following polynomial approximations to the Nevanlinna functions.
 
\begin{prop}\cite[Proposition 5.24]{Sch}\label{thm:An-Dn} For $u,v\in\C$ and $n\ge 0$ we have
\begin{eqnarray*}
A_n(u,v)&:=&(u-v)\sum_{k=0}^nq_k(u)q_k(v)=
a_n\left|\begin{array}{cc}
 q_{n+1}(u)&\;q_{n+1}(v)\\q_{n}(u)&\;q_{n}(v)\end{array}\right|,\\
B_n(u,v)&:=&-1+(u-v)\sum_{k=0}^n p_k(u)q_k(v)=
a_n\left|\begin{array}{cc}
p_{n+1}(u)&\;q_{n+1}(v)\\p_{n}(u)&\;q_{n}(v)\end{array}\right|,\\
C_n(u,v)&:=&1+(u-v)\sum_{k=0}^n q_k(u)p_k(v)=
a_n\left|\begin{array}{cc}
q_{n+1}(u)&\;p_{n+1}(v)\\
q_{n}(u)&\;p_{n}(v)\end{array}\right|,\\
D_n(u,v)&:=&(u-v)\sum_{k=0}^np_k(u)p_k(v)=
a_n\left|\begin{array}{cc}
p_{n+1(}u)&\;p_{n+1}(v)\\
p_{n}(u)&\;p_{n}(v)\end{array}\right|.
\end{eqnarray*}
\end{prop}

The Jacobi operator $(T,D(T))$ has deficiency indices $(1,1)$ and the self-adjoint extensions in
$\ell^2$ can be parametrized as the operators $T_t, t\in\R^*=\R\cup \{\infty\}$ with domain  
\begin{equation}\label{eq:domTt}
D(T_t)=D(T)\oplus \C  (\mathfrak{q}_0 + t\mathfrak{p}_0) \;\mbox{for}\;t\in\R,\quad
D(T_\infty)=D(T)\oplus \C \mathfrak{p}_0
\end{equation}
and defined by the restriction of $T^*$ to the domain, cf. \cite[Theorem 6.23]{Sch}.
We recall that $\mathfrak{p}_0, \mathfrak{q}_0$ are defined in \eqref{eq:frak}.

For $t\in\R^*$ we define the solutions to the moment sequence \eqref{eq:mom}
\begin{equation}\label{eq:Next}
\mu_t(\cdot):=\langle E_t(\cdot)e_0,e_0\rangle,
\end{equation}
where $E_t(\cdot)$ is the spectral measure of the self-adjoint operator $T_t$.

The measures $\mu_t, t\in\R^*$ are precisely those measures $\mu\in V$  for which the polynomials $\C[x]$ are dense in $L^2(\mu)$ according to a famous theorem of M. Riesz, cf. \cite{Ri}.
 They are called N-extremal  in \cite{Ak} and von Neumann solutions in \cite{S}, and they form a compact subset of the set $\ext(V)$ of extreme points of the convex set $V$. However, $\ext(V)$ is known to be a dense subset of $V$.  The N-extremal measures are characterized by the formula
\begin{equation}\label{eq:Npar}
\int\frac{d\mu_t(x)}{x-z}=-\frac{A(z)+tC(z)}{B(z)+tD(z)},\quad z\in \C\setminus\R, t\in\R^*,
\end{equation}
where $A, B, C, D$ are the entire functions given in \eqref{eq:A-D}, cf. \cite[Theorem 7.6]{Sch}. Recall that \eqref{eq:det} holds, so the right-hand side of \eqref{eq:Npar} is a M\"{o}bius transformation in $t$. We note in passing that the solutions $\mu\in V$ different from the N-extremal ones are given in \eqref{eq:Npar}, when $t$ is replaced by a non-degenerate Pick function $\varphi:\C\setminus\R \to\C$, cf. Theorem 7.13 in \cite{Sch}.

We summarize some of the properties of $\mu_t$, which can be found in \cite{Ak} and \cite{Sch}.

\begin{prop}\label{thm:supportNext}
\begin{enumerate}
\item[(i)] The  solution $\mu_t$ is a discrete measure with support equal to the countable zero set $\Lambda_t$ of the entire function $B(z)+tD(z)$, with the convention that $\Lambda_\infty$ is the zero set of $D$. We have $\Lambda_t\subset\R$ for $t\in\R^*$ and
the mass of $\mu_t$ at $u\in\Lambda_t$ equals $||\mathfrak{p}_u||^{-2}$. 
\item[(ii)] The supports of two different N-extremal solutions are disjoint and interlacing. Each point $x_0\in\R$ belongs to the support of a unique N-extremal measure $\mu_t$, where
$t\in\R^*$ is given  as $t=-B(x_0)/D(x_0)$ if $D(x_0)\neq 0$ and $t=\infty$ if $D(x_0)=0$.
\item[(iii)] If $x_0\in\R$ belongs to the support of the N-extremal measure $\mu_t$, then the measure $\mu_t-\mu_t(x_0)\delta_{x_0}$ is determinate.
\end{enumerate}
\end{prop}

Putting $z=0$ in \eqref{eq:Npar}, which is possible when $0\notin \supp(\mu_t)$, leads to
\begin{equation}\label{eq:Npar0}
\int\frac{d\mu_t(x)}{x}=t,\quad t\in\R,
\end{equation}
because $A(0)=D(0)=0, B(0)=-1, C(0)=1$.

The following Proposition shows how to construct an orthonormal basis for $\ell^2$ from the support of an N-extremal measure. We give two proofs, different in style. 
 
\begin{prop}\label{thm:ONB} Let $(u_n)_{n\ge 1}$ be the points of the countable set $\Lambda_t$ supporting $\mu_t$, where $t\in\R^*$. The vectors $\mathfrak{p}_{u_n}/||\mathfrak{p}_{u_n}||, n\ge 1$ form an orthonormal basis for $\ell^2$.
\end{prop}

\begin{proof} First proof. We know that $B(u_n)+tD(u_n)=0$ for $n\ge 1$ if $t\in\R$ and that $D(u_n)=0$ for $n\ge 1$ if $t=\infty$. We then get the orthogonality of the sequence $(\mathfrak{p}_{u_n})$ by
\begin{eqnarray*}
&&(u_n-u_m)\langle \mathfrak{p}_{u_n}, \mathfrak{p}_{u_m}\rangle=D(u_n, u_m)
=B(u_n)D(u_m)-D(u_n)B(u_m)\\
&=&\left\{\begin{array}{ll}
 -tD(u_n)D(u_m)-D(u_n)(-tD(u_m))=0 & \mbox{if} \;\; t\in\R,\\
 0 & \mbox{if}\;\; t=\infty,
 \end{array}
  \right.
\end{eqnarray*}
cf. formula (5.8) in \cite{B:S5}.

Assume next that $c\in\ell^2$ satisfies $\langle c, \mathfrak{p}_{u_n}\rangle=0$ for $n\ge 1$. Then the function $F_c$ from \eqref{eq:FGc} vanishes on $\supp(\mu_t)$ and by \eqref{eq:Pars} we get $c=0$, showing that  $\spn\{\mathfrak{p}_{u_n}\mid n\ge 1\}$ is  dense in $\ell^2$.

Second proof. The mapping $U:L^2(\mu_t)\to \ell^2$ defined by
$$
Uf=\left(\langle f, p_k \rangle_{L^2(\mu_t)}\right)_{k\ge 0}
$$
is a linear isometry by the Parseval identity. As $Up_n=e_n$ the operator  $U$ is unitary.
The functions $||\mathfrak{p}_{u_n}||\delta_{u_n}$ ($\delta_{u_n}$ being 1 on $u_n$ and  0 on $u_m, m\neq n$) form an orthonormal basis in $L^2(\mu_t)$. Therefore
$$
U(||\mathfrak{p}_{u_n}||\delta_{u_n})=\mathfrak{p}_{u_n}/||\mathfrak{p}_{u_n}||
$$
form an orthonormal basis for $\ell^2$.
\end{proof}

 
\section{Proof of Theorem~\ref{thm:concrete}}
 In the proof of Theorem~\ref{thm:concrete} we need the following result about the functions $B(u,v), D(u,v)$. It is  of independent interest.

\begin{prop}\label{thm:1-2} (i) For $\mu\in V, v\in\C$ and any polynomial $p$ the functions $p(u)B(u,v), p(u)D(u,v)$ belong to $L^1(\mu)$ as functions of $u$ and
$$
\int p(u)B(u,v)\,d\mu(u)=\int p(u)D(u,v)\,d\mu(u)=0.
$$ 

(ii) For $v_0\in\R$ we have $B(u,v_0)\notin L^2(\mu)$ for any N-extremal measure $\mu$ different from the N-extremal measure $\mu_s$ supported by $\{u\in\R \mid B(u,v_0)=0\}$, cf. \eqref{eq:N-ext2}, while $B(u,v_0)$ is the zero element in $L^2(\mu_s)$.

(iii) For $v_0\in\R$ we have $D(u,v_0)\notin L^2(\mu)$ for any N-extremal measure $\mu\neq \mu[v_0]$, while $D(u,v_0)$ is the zero element in $L^2(\mu[v_0])$.
 \end{prop} 

\begin{proof} (i). With the notation of \eqref{eq:FGc} we have
\begin{eqnarray*}
p(u)B(u,v)&=&-p(u)+p(u)(u-v)F_{(q_k(v))}(u),\\
 p(u)D(u,v)&=&p(u)(u-v)F_{(p_k(v))}(u),
\end{eqnarray*}
so it is clear from the Cauchy-Schwarz inequality that $p(u)B(u,v), p(u)D(u,v)$ belong to $L^1(\mu)$ as functions of $u$ and furthermore that
$p(u)B_n(u,v)\to p(u)B(u,v)$ and $p(u)D_n(u,v)\to p(u)D(u,v)$  in $L^1(\mu)$ for $n\to\infty$. However,
\begin{equation}\label{eq:int0}
\int p(u)B_n(u,v)\,d\mu(u)=\int p(u)D_n(u,v)\,d\mu(u)=0
\end{equation}
 for $n>\deg(p)$ by Proposition~\ref{thm:An-Dn}, and (i) follows.

(ii). If $B(u,v_0)\in L^2(\mu)$ for an N-extremal measure $\mu$ it follows from \eqref{eq:int0} that $\int p(u)B(u,v_0)\,d\mu(u)=0$ for all polynomials $p$. Since the polynomials are dense in $L^2(\mu)$ we conclude that $B(u,v_0)=0$ for all $u\in\supp(\mu)$. Therefore $\mu$ must be the N-extremal measure supported by $\{u\in\R \mid B(u,v_0)=0\}$.

(iii). If $D(u,v_0)\in L^2(\mu)$ for an N-extremal measure $\mu$, we get that $D(u,v_0)=0$ for all $u\in\supp(\mu)$. This is not possible if $\mu\neq \mu[v_0]$ because then  
$\supp(\mu)$ is disjoint from $\supp(\mu[v_0])$ and this contradicts \eqref{eq:supp}, which also shows that  $D(u,v_0)$ is the zero element in $L^2(\mu[v_0])$.
\end{proof}

\begin{cor}\label{thm:limit} For any $v_0\in\R$ we have
\begin{eqnarray*}
&\lim_{n\to\infty} a_n^2[p_n(v_0)^2+p_{n+1}(v_0)^2]=\infty,\\
&\lim_{n\to\infty} a_n^2[q_n(v_0)^2+q_{n+1}(v_0)^2]=\infty,
\end{eqnarray*}
where $a_n$ is the entry $J_{n+1,n}$  of the Jacobi matrix \eqref{eq:Jac}.
\end{cor}

\begin{proof} To prove the first equation assume by contradiction that there is a subsequence $(n_k)$ such that $a_{n_k}^2[p_{n_k}(v_0)^2+p_{n_k+1}(v_0)^2]$ is bounded, say by $M$.  We know that $D_n(u,v_0)$ from Proposition~\ref{thm:An-Dn} tends locally uniformly in $u$ to $D(u,v_0)$ for $n\to\infty$, and hence for any $R>0$ and any $\mu\in V$ 
\begin{eqnarray*}
\int_{-R}^R D(u,v_0)^2\,d\mu(u)&=&\lim_{k\to\infty}\int_{-R}^R D_{n_k}(u,v_0)^2\,d\mu(u)\\
&\le&\limsup_{k\to\infty}\int_{-\infty}^\infty D_{n_k}(u,v_0)^2\,d\mu(u)\\
&=&\limsup_{k\to\infty} a_{n_k}^2[p_{n_k}(v_0)^2+p_{n_k+1}(v_0)^2]\le M.
\end{eqnarray*}
As $R>0$ is arbitrary, $D(u,v_0)$ belongs to $L^2(\mu)$ for any $\mu\in V$, but this is not true by Proposition~\ref{thm:1-2} (iii). 

The proof of the second equation is similar.
\end{proof}

The following Lemma will be crucial for the proof of Theorem~\ref{thm:concrete}.

\begin{lem}\label{thm:Hilbert} Let $(x_n)_{n\ge 1}$ denote an orthonormal basis in a Hilbert space $\mathcal H$, let $(a_n)_{n\ge 2}$ be a sequence of complex numbers and $\xi\in\mathcal H$ a vector satisfying $\langle \xi, x_1\rangle\neq 0.$

Then $\spn\{x_n+a_n\xi\mid n\ge 2\}$ is dense in $\mathcal H$ iff $\sum_{n=2}^\infty
|a_n|^2=\infty$.

Moreover, $\spn\{x_n+a_n\xi\mid n\ge 3\}$ is never dense in $\mathcal H$. 
\end{lem}

\begin{proof} 1. Assume first the existence of a vector $w\in\mathcal H, w\neq 0$ such that
$\langle w, x_n+a_n\xi\rangle=0$ for $n\ge 2$. Then we have
$\langle w,x_n\rangle=-\overline{a_n}\langle w,\xi\rangle$ for $n\ge 2$.
 
We claim that  $\langle w, \xi\rangle\neq 0$, for if $\langle w, \xi\rangle=0$, we get $\langle w, x_n\rangle=0$ for $n\ge 2$ so $w=\langle w, x_1\rangle x_1$, and hence
$$
0=\langle w, \xi\rangle=\langle w, x_1\rangle \langle x_1, \xi\rangle.
$$
By assumption $\langle \xi, x_1\rangle\neq 0$, so $\langle w, x_1\rangle =0$ and together with the equations $\langle w, x_n\rangle=0$ for $n\ge 2$, we get $w=0$, which is a contradiction. 

Using $\langle w, \xi\rangle\neq 0$, we get by Bessel's inequality
$$
||w||^2\ge \sum_{n=2}^\infty |\langle w, x_n\rangle|^2=|\langle w,\xi\rangle|^2\sum_{n=2}^\infty|a_n|^2,
$$
showing that $\sum_{n=2}^\infty|a_n|^2<\infty$.

2. Assume next that $\sum_{n=2}^\infty|a_n|^2<\infty$, and define $w\in\mathcal H$ by
$$
w=a_1x_1-\sum_{n=2}^\infty \overline{a_n}x_n
$$
with $a_1$ such that $\langle w, \xi\rangle=1$, i.e., such that
$$
a_1\langle x_1,\xi\rangle-\sum_{n=2}^\infty \overline{a_n}\langle x_n, \xi\rangle =1,
$$
which is possible as $\langle \xi, x_1\rangle\neq 0$.

For $n\ge 2$ we find
$\langle w, x_n+a_n\xi\rangle=-\overline{a_n}+\overline{a_n}\langle w, \xi\rangle=0$.
Note that  $w\neq 0$ because either $a_1\neq 0$ or if $a_1=0$, i.e., $\sum_{n=2}^\infty \overline{a_n}\langle x_n, \xi\rangle =-1$, then $a_n\neq 0$ for some $n\ge 2$.

3. To see that $\spn\{x_n+a_n\xi \mid n\ge 3\}$ is not dense in $\mathcal H$ we define
$y:=\langle x_1, \xi\rangle x_2-\langle x_2, \xi\rangle x_1$ and note that $y\neq 0$.
For $n\ge 3$ we find
\begin{eqnarray*}
&&\langle y, x_n+a_n\xi\rangle=\langle \langle x_1,\xi\rangle x_2-\langle x_2, \xi\rangle x_1, a_n\xi\rangle\\
&=& \overline{a_n}\left(\langle x_1, \xi\rangle \langle x_2, \xi\rangle-\langle x_2, \xi\rangle\langle x_1, \xi\rangle\right)=0.
\end{eqnarray*}
\end{proof}

\noindent{\it Proof of Theorem~\ref{thm:concrete} and Remarks~\ref{thm:optimal1} and \ref{thm:optimal2}} 

(i). For $v_0\in\R$ we shall see that the density of $P(v_0)$ in $\ell^2$ follows from
 Lemma~\ref{thm:Hilbert}. We write the numbers in $\{u\in\R \mid D(u,v_0)=0\}$ as a sequence $(u_n)_{n\ge 1}$ with $u_1=v_0$. According to Proposition~\ref{thm:ONB} the vectors $x_n=\mathfrak{p}_{u_n}/||\mathfrak{p}_{u_n}||, n\ge 1$ form an orthonormal basis for $\ell^2$. For $\xi=\mathfrak{p}_{v_0}$ we have $\langle \xi, x_1\rangle=||\mathfrak{p}_{v_0}||\neq 0$ and
$$
P(v_0)=\spn\{x_n+a_n\xi \mid n\ge 2\},\quad a_n=B(u_n,v_0)/||\mathfrak{p}_{u_n}||.
$$
Using Proposition~\ref{thm:supportNext} (i) we get
$$
\sum_{n=2}^\infty \frac{B(u_n,v_0)^2}{||\mathfrak{p}_{u_n}||^2}=\int_{\R\setminus\{v_0\}}B(u,v_0)^2\,d\mu[v_0](u),
$$  
which is infinite by Proposition~\ref{thm:1-2} (ii).

The subspace $P'(v_0)$ equals $\spn\{x_n+a_n\xi \mid n\ge 3\}$, so it is not dense in $\ell^2$ by Lemma~\ref{thm:Hilbert}.

(ii).  The case about $Q(v_0)$  can be deduced from case (i) by using the observation that
the polynomials $(q_{n+1}(x)/q_1(x))_{n\ge 0}$ are the orthonormal polynomials associated with the truncated Jacobi matrix $J^{(1)}$ obtained from $J$ by removing the first row and column. See \cite[p. 28]{Ak}, \cite[p. 131, p. 144]{B:S5} and \cite{Pe} for details.

(iii). We finally consider $M(v_0)$. From \eqref{eq:N-ext2}  we have $B(u,v_0)=0$ if and only if $u\in\supp(\mu_s)$, where 
$s=-A(v_0)/C(v_0)$, interpreted as $s=\infty$ if $C(v_0)=0$.
These zeros form a sequence denoted $(u_n)_{n\ge 1}$. The vectors $x_n=\mathfrak{p}_{u_n}/||\mathfrak{p}_{u_n}||, n\ge 1$ form an orthonormal basis for $\ell^2$, and with $\xi=\mathfrak{q}_{v_0}$ and $a_n=-D(u_n,v_0)/||\mathfrak{p}_{u_n}||, n\ge 1$  we have
$$
\langle \xi, x_1\rangle=1/((u_1-v_0)||\mathfrak{p}_{u_1}||)\neq 0,\quad M(v_0)=\spn\{x_n+a_n\xi \mid n\ge 1\}.
$$  
Furthermore,
$$
\sum_{n=1}^\infty a_n^2= \int D(u,v_0)^2\,d\mu_s(u)=\infty
$$
by Proposition~\ref{thm:1-2}(iii) since $\mu_s\neq \mu[v_0]$ because $B(v_0,v_0)=-1$.

It is clear that also $\sum_{n=2}^\infty a_n^2=\infty$, so by Lemma~\ref{thm:Hilbert} the subspace
$M'(v_0)=\spn\{x_n+a_n\xi \mid n\ge 2\}$ of $M(v_0)$ is dense in $\ell^2$, while
$M''(v_0)=\spn\{x_n+a_n\xi \mid n\ge 3\}$ is not dense in $\ell^2$. \hfill$\square$

\medskip
{\bf Acknowledgment} The authors want to thank two  independent referees for their careful comments to the manuscript.  This has led to improvements of the paper.


\begin{thebibliography}{12}
\bibitem{Ak}  N.~I. Akhiezer,  {\it The Classical Moment Problem and Some Related
 Questions in Analysis}. English translation, Oliver and Boyd, Edinburgh,
 1965.

\bibitem{B:C} C.~Berg and J.~P.~R.~Christensen,  {\it Density questions in the classical theory of moments}.  Ann. Inst. Fourier  {\bf 31}, no. 3 (1981), 99--114.

\bibitem{B:S1} C.~Berg and R.~Szwarc,  {\it The Smallest Eigenvalue of Hankel Matrices}. Constr. Approx. {\bf 34} (2011), 107--133.

\bibitem{B:S2} C.~Berg and R.~Szwarc,  {\it Inverse of infinite Hankel moment matrices}. SIGMA, Symmetry Integrability Geom. Methods Appl. {\bf 14} (2018), paper 109, 48 pages.

\bibitem{B:S3} C.~Berg and R.~Szwarc,  {\it Closable Hankel Operators and Moment Problems}.
Integr. Equ. Oper. Theory {\bf 92}(1) (2020), 1--9.

\bibitem{B:S4} C.~Berg and R.~Szwarc, {\it Self-adjoint operators associated with Hankel moment matrices}. J. Funct. Anal., {\bf 283} (2022), 109674. 

\bibitem{B:S5} C.~Berg and R.~Szwarc, {\it Indeterminate Jacobi operators}. J. Operator Theory {\bf 93:1} (2025), 123--145.

\bibitem{Pe} H.~L.~Pedersen, {\it The Nevanlinna matrix of entire functions associated with a shifted indeterminate Hamburger moment problem}. Math. Scand. {\bf 74} (1994), 152--160.

\bibitem{Ri} M.~Riesz, {\it Sur le probl{\`e}me des moments et le th{\'e}or{\`e}me de Parseval correspondant}. Acta Litt. Ac. Sci. Szeged {\bf 1} (1923), 209--225.

\bibitem{S} B.~Simon, {\it The classical moment problem as a self-adjoint finite difference operator}. Adv. Math. {\bf 137} (1998), 82--203.

\bibitem{Sch} K.~Schm\"{u}dgen, {\it The Moment Problem}. Graduate Texts in Mathematics Vol. 277. Springer International Publishing AG 2017. 

\end{thebibliography}

\noindent
Christian Berg\\
Department of Mathematical Sciences, University of Copenhagen\\
Universitetsparken 5, DK-2100 Copenhagen, Denmark\\
e-mail: {\tt{berg@math.ku.dk}}

\vspace{0.4cm}
\noindent
Ryszard Szwarc\\
Institute of Mathematics, University of Wroc{\l}aw\\
pl.\ Grunwaldzki 2/4, 50-384 Wroc{\l}aw, Poland\\ 
e-mail: {\tt{szwarc2@gmail.com}}

\end{document}





